% 本文件由 LaTeX 排版 Agent 根据有效 translation.json 自动生成。
% author=Theodoret; work_id=2701; title=驳斥济利禄十二绝罚之辩词

\chapter{驳斥济利禄\WorkTitle{十二绝罚}之辩词}

% structure: p0001 source=h1 target=omitted reason=page-title-already-rendered-by-parent

（舒尔茨编\WorkTitle{作品集}卷I，页1起；米涅\WorkTitle{拉丁教父集}卷76，栏391）\par

\Emph{驳第一条}——然而我们这些遵循福音作者的话语的人声明：天主圣言并非因本性而成为血肉，也未曾变为血肉；因为天主性是不变不更的。为此，达味先知也说：\BibleQuote{祢却常存，祢的岁月没有穷尽。}伟大的保禄，真理的宣讲者，在\BibleBook{希伯来}{书}中指出这句话是针对圣子说的。又在另一处，天主藉先知说：\BibleQuote{我是上主，从不改变。}既然天主性是不变不更的，它就不可能改变或变动。如果那不变者不能改变，那么天主圣言成为血肉就不是通过变化，而是取了血肉，并住在我们中间，一如福音作者所言。神圣的保禄在\BibleBook{斐理伯}{书}中清楚地表达了这一点，他写道：\BibleQuote{你们该怀有基督耶稣所怀有的心情：祂虽具有天主的形体，却没有以自己与天主同等为应当把持不舍的，反而空虚自己，取了奴仆的形体。}从这些话显然可见，天主的形体并未改变为奴仆的形体，而是保持其本有，取了奴仆的形体。所以，天主圣言并非成了血肉，而是取了有理性的、生活的血肉。祂自己并非由童贞女自然地受孕、被塑造、被形成，并从她获得存在的开端；祂是万世之前的天主，与天主同在，与圣父同在，被认识、被朝拜；祂却在自己的母腹中为自己塑造了一座圣殿，并与那被塑造和被生育者同在。因此，我们也称那神圣童贞女为\OriginalTerm{天主之母}{\GreekText{θεοτόκος}}，不是因为她以自然方式生了天主，而是因为她生了与那塑造祂的天主结合的人。此外，如果那在童贞女腹中被塑造的不是人，而是万世之前的天主圣言，那么天主圣言就是圣神的受造物了。因为加俾额尔说：\BibleQuote{那在她内受生的，是出于圣神。}但如果天主的独生圣言是非受造的，与父同体、同永恒，那么祂就不再是圣神所塑造或创造的。如果圣神并未在童贞女腹中塑造天主圣言，那么我们就应理解为：奴仆的形体被塑造、被形成、被孕生、被诞生。但由于那形体并没有脱离天主的形体，而是作为一座圣殿，容纳了居住其中的天主圣言，正如保禄所说：\BibleQuote{因为天主乐意叫整个的圆满居住在祂内}，\BibleQuote{具有形体地}，所以我们称童贞女不是\OriginalTerm{人之母}{\GreekText{ἀνθρωποτόκος}}，而是\OriginalTerm{天主之母}{\GreekText{θεοτόκος}}；前者指塑造和孕生，后者指结合。为此，那诞生的婴孩被称为厄玛奴耳，既非与人性分离的天主，也非剥夺了天主性的人。因为厄玛奴耳被解释为\BibleQuote{天主与我们同在}，一如福音所言；而\BibleQuote{天主与我们同在}这句话，既显明了那为我们而从我们中被取的那位，又宣告了那取了祂的天主圣言。因此，婴孩因那取人的天主而被称为厄玛奴耳，童贞女因天主的形体与所取的奴仆形体的结合而被称为\OriginalTerm{天主之母}{\GreekText{θεοτόκος}}。因为天主圣言并未变成血肉，而是天主的形体取了奴仆的形体。\par

\Emph{驳第二条}——我们遵从宗徒的神圣教导，宣认一位基督；并因着结合，我们称同一基督既是天主又是人。但我们完全不知道所谓\Quote{按位格的结合}，因为它对圣经和解释圣经的教父来说都是陌生而外来的。如果这些声明的作者所谓\Quote{按位格的结合}是指血肉与天主性的混合，我们将全力反对他的说法，驳斥他的亵渎；因为混合必然导致混淆，而承认混淆便摧毁了各本性的独特性。正在混合的事物不会保持其所是，若在天主圣言和达味的后裔身上断言这一点，是极其荒谬的。我们必须听从主，祂显示了两个本性，并对犹太人说：\BibleQuote{你们拆毁这座圣殿，三天之内我要把它重建起来。}但如果有了混合，那么天主就不再是天主，圣殿也不再被认作圣殿；那么圣殿就是天主，天主就是圣殿。这就是混合论所涉及的结果。而主对犹太人说\BibleQuote{你们拆毁这座圣殿，三天之内我要把它重建起来}就完全是多余的了。祂本应说：你们拆毁我，三天之内我要被重建起来——如果真的有任何混合和混淆的话。但事实是，祂显明圣殿遭受拆毁，而天主使它复活。因此，我认为他们用来代替混合的\Quote{按位格的结合}是多余的。只要提及\Quote{结合}就足够了，它既显明本性的特质，又教导我们朝拜唯一的基督。\par

\Emph{驳斥第三条。}——所用术语之意涵模糊晦涩。谁需要被告知，联合与结合并无区别？结合是分离部分的结合，而联合是区分部分的联合。这些措辞的绝顶聪明之作者将本来一致的事物说成似乎不一致。他说，以联合来结合位格是错误的，而应以结合，且是自然的结合来结合。或许他说这话时不知道自己在说什么；若他知道，便是亵渎。自然具有强制性且非自愿；例如，若我说我们自然地饥饿，则我们并非出于自由意志而是出于必然感到饥饿；的确，若贫苦之人有能力凭自己意志停止饥饿，他们早就不再乞讨了；我们自然地口渴；我们自然地睡眠；我们自然地呼吸；我重申，所有这些行动都属于非自愿的范畴，而一个不能再进行这些行动的人必然停止存在。因此，若天主性与奴仆性在结合中的结合是自然的，那么天主圣言与奴仆性的结合便是出于必然的强制，而非因祂施以慈爱，而宇宙的立法者将被发现是必然法则的追随者。有福的保禄并未如此教导我们；相反，我们被教导祂取了奴仆性，并\Quote{空虚了自己}\BibleRef{斐 2:7}；而\Quote{空虚了自己}这一表达表明了自愿的行动。因此，若祂是出于旨意和意志而与从我们取的人性结合，那么\Quote{自然}一词的添加便是多余的。承认结合就够了，而结合是就区分的事物而言的，因为若没有区分，则结合绝不能为人所知。因此，对结合的理解意味着对区分的先行理解。那么他怎能说位格或性不应被区分？他始终知道天主圣言的位格在万世之前已是完美的；祂所取的奴仆性也是完美的；正因如此他说\Quote{位格}而非\Quote{位格}。因此，若每一性都是完美的，且两者结合，显然在形式天主取了形式奴仆之后，虔敬迫使我们承认一个子和基督；而将结合的位格或性说成两个，非但不荒谬，反而符合必然之理。因为即使在一个人身上，我们也区分性，称可朽的性为身体，不朽的性为灵魂，两者都是人；那么承认取者天主与被取者人的特性，就更合乎正确理性了。我们发现蒙福的保禄将一个人分为两个，在一处说：\Quote{我们外在的人虽在毁坏，内在的人却日日更新}\BibleRef{格后 4:16}，在另一处说：\Quote{因为我按内在的人喜悦天主的法律}\BibleRef{罗 7:22}，又说：\Quote{使基督因信住在你们心中}\BibleRef{弗 3:17}。既然宗徒将同时之性的自然结合加以区分，那么这位用其他术语向我们描述混合的人，有何理由在我们将永远天主与被取于末期的人的性的特性加以区分时，指控我们不虔敬呢？\par

\Emph{驳斥第四条}——这些话也与前面的话类似。他假定已有混合，然后声称在圣福音和宗徒著作中所用的术语有区别。他讲这些话时自夸他同时与\Person{亚略}、\Person{欧诺米}及其他异端首领作战。那么，让这位精确的神学教授告诉我们，他如何驳斥异端者的亵渎，同时将\OriginalTerm{仆人的形像}{form of the servant}所谦卑而适当地讲出的话应用于\OriginalTerm{天主圣言}{God the Word}。他们这样做，确实是主张天主子是低劣的、受造物、被造品和仆人。那么，我们这些持相反意见、明认圣子与天主圣父\OriginalTerm{同一性体}{one substance}、\OriginalTerm{同永恒}{co-eternal}、是宇宙的创造者、造物主、美化者、治理者、管理者、全智者、全能者，或更确切地说，祂本身就是能力、生命和智慧的人，应该把以下的话归给谁呢？\BibleQuote{我的天主，我的天主，祢为什么舍弃了我？}或\BibleQuote{父啊！若是可能，就让这杯离开我吧！}或\BibleQuote{父啊，救我脱离这时辰}或\BibleQuote{那时辰没有人知道，连人子也不知道}，以及所有由祂和圣宗徒们关于祂所讲的写下的谦卑的话。我们把疲倦和睡眠归给谁？把\OriginalTerm{无知}{ignorance}和恐惧归给谁？谁需要天使的帮助？如果这些属于天主圣言，那么智慧怎么会无知？如果智慧受无知影响，怎么能被称为智慧？祂说凡父所有的都是祂的，既然没有父的知识，怎能说这话？因为祂说：\BibleQuote{只有父知道那一天}。如果祂不是完全拥有生者所有的一切，怎能是祂的不可改变的\OriginalTerm{肖像}{image}？那么，如果祂说祂无知时讲的是真话，任何人都可能这样想祂。但如果祂知道那一天，却为了隐藏而说祂无知，你就会看到结论是多么亵渎。因为真理说谎，如果它有与真理相反的性质，就不能恰当地被称为真理。但如果真理不说谎，那么天主圣言对自己所创造并且自己已定的、决定审判世界的那一天就不是无知的，而是拥有父的知识，作为不改变的肖像。那么，无知不属于天主圣言，而属于仆人的形像，那形像当时只知道内住的神性所启示的程度。同样的立场可以用于其他类似情况。例如，天主圣言对父说\BibleQuote{父啊！若是可能，就让这杯离开我吧！但不要照我的意愿，而要照祢的意愿}，这怎么合理呢？由此必然产生的荒谬之处不少。首先，父和子不同心，父愿意一件事，子愿意另一件事，因为祂说：\BibleQuote{但不要照我的意愿，而要照祢的意愿}。其次，我们将看见子有极大的无知，因为祂将不知道那杯能否离开祂；但这对天主圣言来说是极大的不虔和亵渎。因为祂准确地知道\OriginalTerm{救恩计划的奥秘}{mystery of the œconomy}，祂为此来到我们中间，甘愿取了我们的本性，并空虚了自己。为此，祂也预言给圣宗徒们：\BibleQuote{看，我们上耶路撒冷去，人子要被交出来……交在外邦人手中，被戏弄、鞭打、钉死，第三天祂要复活。}那么，预言这些事、并在\Person{伯多禄}劝阻时责备他的那一位，怎能自己又劝阻这些事发生，明明知道一切将要发生的事呢？\Person{亚巴郎}在许多世代以前就看见祂的日子而欢喜，同样\Person{依撒意亚}、\Person{耶肋米亚}、\Person{达尼尔}、\Person{匝加利亚}和众先知团体都预言了祂的救恩性苦难，而祂自己却无知，请求解脱并劝阻，尽管这事注定要为世界的救恩而发生，这难道不荒谬吗？因此，这些话不是天主圣言的话，而是仆人的形像的话，它害怕死亡，因为死亡尚未被消灭。确实，天主圣言容许这些表达，留出恐惧的空间，以便成为人的那一位的本性得以显明，并防止我们认为亚巴郎和达味之子是不真实的或只是表象。不虔的异端者团伙因为持有这些意见而产生了这种亵渎。因此，我们将神圣地所说和所做的一切应用于天主圣言；另一方面，在谦卑中所说所做的一切，我们将它与仆人的形像联系起来，免得我们染上\Person{亚略}和\Person{欧诺米}的亵渎。\par

\Emph{驳斥第五条}——我们声明，天主圣言像我们一样分享了血肉和人的不朽灵魂，这是由于与它们\OriginalTerm{结合}{union}的缘故；但我们不仅拒绝说天主圣言因任何变化而成了血肉，而且斥责这样做的人为不虔，并且可见这与所提出的措辞相反。因为如果圣言改变成了血肉，祂就没有与我们分享血肉；但如果祂分享了血肉，祂就是作为与血肉不同的另一位而分享；如果血肉是不同于祂的某物，那么祂就没有改变成血肉。因此，我们使用\Quote{分享}一词时，将那位取了和被取的一位作为一位圣子来敬拜。但我们承认\OriginalTerm{性体的区分}{distinction of natures}。我们不反对\OriginalTerm{带天主的人}{man bearing God}这一术语，许多圣教父用过它，其中一位是\Person{大巴西略}，他在关于圣神的致\Person{盎斐罗基}的论述以及《第五十九篇圣咏》的解释中使用了这个术语。但我们称祂为\Quote{带天主的人}，不是因为祂领受了某种特殊的神恩，而是因为祂拥有与圣子结合的\OriginalTerm{全部天主性}{all the Godhead of the Son}。因为蒙福的\Person{保禄}在他的解释中说：\BibleQuote{你们要谨慎，恐怕有人用哲学和虚妄的欺骗，按照人的传统，按照世俗的原理，而不是按照基督，把你们掳去。因为在天主性内，一切圆满都有形有实地居住在祂内。}\par

\Emph{驳斥第六条。}——蒙福的保禄称天主圣言所取的那一位为\BibleQuote{奴仆的形体}，但由于那取用是在结合之前，而蒙福的保禄在论及取用时称所取的本性为\BibleQuote{奴仆的形体}，在结合完成后，\BibleQuote{奴役}之名不再有位置。因为看到宗徒写信给信他的人说：\BibleQuote{所以你不再是奴隶，而是儿子}，而主对他的门徒说：\BibleQuote{从今以后，我不再称你们为仆人，反而称你们为朋友}；更何况我们本性的初果，借着祂我们甚至获得了义子恩赐的奖赏，将脱离仆人的称号。因此，我们甚至承认\BibleQuote{奴仆的形体}也是天主，因为\OriginalTerm{天主的形体}{form of God}与它结合；我们听凭先知的话，当他称那婴孩也\OriginalTerm{厄玛奴耳}{Emmanuel}，以及那出生的孩子为\BibleQuote{伟大谋略的天使、奇妙的策士、强有力的神、和平的君、永恒之父}。然而，同一位先知，即使在结合之后，在宣告所取的本性时，称那出于亚巴郎后裔者为\BibleQuote{仆人}，话是：\BibleQuote{你是我的仆人以色列，我要在你身上受光荣}；又说：\BibleQuote{那从母腹形成我作他仆人的上主这样说}；稍后又说：\BibleQuote{看，我立你作人民的盟约、外邦人的光明，使你成为我的救恩直到地极}。但从母腹形成的那一位不是天主圣言，而是奴仆的形体。因为天主圣言成为血肉并非通过变化，而是取了有灵魂的血肉。\par

\Emph{驳斥第七条。}——如果人的本性是会死的，而天主圣言是生命和生命的赐予者，且复活了被犹太人毁坏的圣殿，并把它带上天堂，那么\OriginalTerm{奴仆的形体}{form of the servant}怎能不透过\OriginalTerm{天主的形体}{form of God}受到光荣呢？因为如果它原本和按其本性是会死的，通过它与天主圣言的结合而成为不死的，那么它就接受了它原先没有的；接受了它原先没有的之后，并被光荣了，它是由赐予者光荣的。因此，宗徒也大声说：\BibleQuote{按照祂强大能力的运作，就是在基督身上所运作的，使祂从死人中复活的能力}。\par

\Emph{驳斥第八条。}——正如我常说的，我们向主基督所献的光荣赞颂是唯一的，我们承认祂同时是天主和人，正如结合的方式教导我们的；但我们不会不愿谈论两性各自的特性。因为天主圣言并未变成血肉，人也没有失去其自身而变成天主的本性。因此，我们朝拜主基督，同时保持两性各自的特性。\par

\Emph{驳第九条}—— 此人在此显然胆敢绝罚不仅今日持虔敬意见之人，亦包括昔日真理之宣讲者；甚至神圣福音的作者们、圣宗徒之团体，以及加俾额尔总领天使。因为他正是那一位，在怀孕之前首先按肉身宣告了基督的诞生；当玛利亚问说：\BibleQuote{这事如何成就？因为我不认识男人}，他回答说：\BibleQuote{圣神要临于你，至高者的能力要庇荫你，因此那要诞生的圣者，将称为天主子}。又对若瑟说：\BibleQuote{不要怕娶你的妻子玛利亚，因为那在她内受孕的是出于圣神}。福音作者也说：\BibleQuote{他的母亲玛利亚许配给若瑟……她因圣神有孕的事已显示出来}。主亲自进入犹太人的会堂，拿起依撒意亚先知书，读了其中一段：\BibleQuote{上主的神临于我，因为他给我傅了油}等话，然后说：\BibleQuote{今天这应验在你们耳中了}。蒙福的伯多禄在向犹太人讲道时说：\BibleQuote{天主怎样以圣神傅了纳匝肋的耶稣}。依撒意亚在很久以前预言说：\BibleQuote{从叶瑟的树干要生出一根枝条，从他的根上要发出一个嫩芽；上主的神要住在他身上，即智慧和聪敏的神，超见和刚毅的神，明达和敬畏上主的神}；又说：\BibleQuote{看哪，我的仆人，我所扶持的，我所拣选的，我心所喜悦的！我已将我的神赐给他，他要向外邦人施行审判}。福音作者也将这见证写入了自己的作品中。主在福音中对犹太人说：\BibleQuote{我若是仗赖天主之神驱魔，那么天主的国已经来到你们中间了}。若翰说：\BibleQuote{那派遣我来以水施洗的曾对我说：你看见圣神降下来，停在谁身上，谁就是用圣神施洗的人}。因此，这位神圣法令的精察者不仅绝罚了先知、宗徒，甚至加俾额尔总领天使，更使他的亵渎直达世界救主本身。因为我们已表明，主在读了\BibleQuote{上主的神临于我，因为他给我傅了油}这段话后，对犹太人说：\BibleQuote{今天这应验在你们耳中了}。他对那些说他是靠贝耳则步驱魔的人回答说，他是靠天主之神驱魔。但我们主张，被圣神所形成和\OriginalTerm{受傅}{anointed}的，并非与圣父同体同永的\OriginalTerm{天主圣言}{God the Word}，而是他在末世所取的人性。如果他说\OriginalTerm{圣子之神}{Spirit of the Son}是出于同一性体并由圣父所发，我们就承认这话，认为符合真宗教。但如果他说圣神是属于圣子，或通过圣子而有根源，我们就拒绝这话，视为亵渎和不敬。因为我们相信主说：\BibleQuote{那由圣父所发的圣神}；同样极其神圣的保禄说：\BibleQuote{我们不是领受了世界的神，而是领受了那由天主而来的神}。\par

\Emph{驳第十条：}——不变的本性并未转变为血肉的本性，而是取了我们的人性，并将其安置在大司祭之上，正如蒙福的保禄所教导的：\Quote{凡从人间选立的大司祭，是为人的事奉天主的事，以便奉献供物和赎罪的牺牲；他能同情无知和迷途的人，因为他自己也为软弱所困。因此，他理当为人民也为自己的罪献祭。}稍后解释这一点时他说：\InnerQuote{正如亚郎，基督也是如此。}然后指出所取人性的软弱，他说：\InnerQuote{祂在血肉之日，以大声哀号和眼泪，向那能救祂脱离死亡的天主献上祈祷和恳求，因祂的虔敬而获得了垂允。祂虽然是儿子，却由所受的苦难学习了服从，且得以成全后，为一切服从祂的人成了永远救恩的根源，被天主称为按照默基瑟德品位的大司祭。}那么，是谁借着德行的劳苦而得以成全，本性却并不完善？是谁借由经验学习了服从，在经验之前对此一无所知？是谁怀着虔敬生活，以大声哀号和眼泪献上恳求，不能救自己，却向那能救祂的天主呼求，恳求脱离死亡？不是天主圣言——那不受苦难、不死、无形体者，祂的记忆是喜乐和脱离眼泪，\InnerQuote{因为祂擦去了各人脸上的眼泪}，先知又说：\InnerQuote{我想起天主，便喜乐。}祂加冕那些虔敬生活的人，\InnerQuote{祂在事未成之前已知道一切}，\InnerQuote{凡圣父所有的，祂都有}；祂是圣父不变的本体之像，\InnerQuote{祂在自身内显示圣父。}相反，那由祂所取的达味后裔，是凡人、可受苦、惧怕死亡的；虽然这本身后来借着与取了它的天主结合，摧毁了死亡的能力；它行完一切义德，对若翰说：\InnerQuote{你暂且容许吧，因为我们应该这样尽全义。}这取了默基瑟德品位的司祭名号，因为它披上了本性的软弱——并非全能的天主圣言。因此，稍早前蒙福的保禄也曾说：\InnerQuote{我们所有的，不是一位不能同情我们软弱的大司祭，而是一位在各方面受过试探，与我们相似，只是没有罪过。}是为我们的缘故而从我们取来的那本性，经历了我们的情感，却没有罪，而非那为我们的救恩取了它的那位。在这部分主题的开端，他教导我们说：\InnerQuote{你们应细察我们宣认的宗徒和大司祭耶稣，祂对委派祂的那位是忠信的，如同梅瑟在祂的全家一样。}但任何持守正确信仰的人，都不会称那非受造的、无始的、与圣父同永恒的天主圣言为受造物；相反，那达味后裔，那无任何罪过而被立为我们的大司祭和祭品者，在将自己为我们的缘故奉献给天主后，祂在自身内结合了圣言——出自天主的圣言——与祂不可分离地相连。\par

\Emph{驳第十一条：}——依我之见，他看似留心真理，以便借真理掩饰其不健全的观点，从而不被发觉与异端者主张同样的教义。但没有什么比真理更强有力，它以其光芒揭露虚假的黑暗。借助其光照，我们将显明他的异端信仰。首先，他从未提及有灵的肉身，也未承认所取的人性是完善的，而是处处依照阿波里纳里的教导谈论肉身。其次，在引入\Quote{混合}这一概念时，他用了其他术语，并将其纳入自己的论证；因为在那处，他明确说主的肉身是没有灵魂的。他说：若有人说主的肉身并非真正属于那出自天主圣父的圣言自己的肉身，而是属于另一位，此人该受绝罚。由此可见，他不承认天主圣言取了灵魂，而只取了肉身，并且圣言本身取代了灵魂在肉身中的位置。我们则主张，主的肉身内含有生命，是赋予生命且有理智的，因为赋予生命的天主性与其结合。而他自己也不情愿地承认两个本性的差异，他谈论肉身和\Quote{天主圣言}，并称其为\Quote{祂自己的肉身}。因此，天主圣言并未转变为血肉的本性，而是拥有祂自己的肉身，即所取的人性，并借着结合使其成为赋予生命的。\par

\Emph{驳第十二条：}——苦难属于可受苦者；不可受苦者超乎苦难。因此，是奴仆的形体受苦，天主的形体当然与之共处，并因苦难所带来的救恩而容许它受苦，且由于结合而将苦难视为己有。所以，并非基督受苦，而是天主从我们中所取的那人受苦。为此，蒙福的依撒意亚在他的预言中宣告：\InnerQuote{忧愁的人，熟悉痛苦。}主基督自己也对犹太人说：\InnerQuote{你们为什么想要杀我，一个给你们讲真理的人？}但面临死亡威胁的不是生命本身，而是拥有可朽坏本性者。在另一处，主对犹太人也教导这一训诲说：\InnerQuote{你们拆毁这座圣殿，三天之内我要把它重建起来。}因此，被拆毁的是（出自达味的圣殿），在它被拆毁后，由那在万世之前无苦难地从圣父所生的独生圣言重建。\par

% structure: p0015 source=h2 target=section reason=relative-heading-rank; base=h2

\section{基利尔反对聂斯多略的绝罚条款}

（曼西《会议集》第四卷，第1067–1082页；米涅《拉丁教父全集》第76卷，第391栏。聂斯多略针对基利尔的绝罚条款见于哈杜因《会议集》第一卷，第1297页。）\par

一、若有人拒绝承认厄玛努耳是真正的天主，因而圣童贞是\OriginalTerm{天主之母}{\GreekText{θεοτόκος}}，因为她按肉体方式生下了降生成人的天主圣言，此人当受绝罚。\par

二、若有人拒绝承认天主圣父的圣言与肉体位格结合，并且是同一基督，兼具天主性和人性，此人当受绝罚。\par

三、若有人将同一基督在结合后分开位格，仅以尊严、独立或特权等关系连接，而不以合乎本性的结合相联，此人当受绝罚。\par

四、若有人将福音作者和宗徒们著作中的言辞，无论是出自圣人对基督所说的，或是基督论及自己的，分开归属于两个位格，并以此把一些话应用于与天主圣言分离的常人，而把另一些话应用于天主圣父的圣言本身，此人当受绝罚。\par

五、若有人胆敢声称基督是带着天主的凡人，而否认祂是真正的天主、唯一的圣子，且由于圣言成为血肉，分享了与我们同样的血肉，此人当受绝罚。\par

六、若有人胆敢声称天主圣父的圣言是基督的天主或主，而不承认同一圣言同时是天主和人，且按圣经所言成为血肉，此人当受绝罚。\par

七、若有人说耶稣作为人被天主圣言推动，并因作为另一者而蒙受独生者的荣耀，此人当受绝罚。\par

八、若有人胆敢主张被提升的人应与天主圣言同受敬拜、同受荣耀、同被称为天主（因为介词\Quote{同}的连续使用强制此涵义），而不以同一敬拜行为尊崇厄玛努耳并以同一颂词赞美祂，因为祂是成为血肉的圣言，此人当受绝罚。\par

九、若有人说主耶稣基督藉圣神得荣耀，将圣神视为外来能力而藉其行动，并因此获得能力制服不洁之灵、在人中行奇迹，而不承认圣神是祂自己的，祂也藉此施行奇迹，此人当受绝罚。\par

十、圣经记载基督是我们所宣认的大司祭和宗徒，祂为了我们将自己献于天主和我们的圣父作为馨香的祭品。因此，若有人声称天主圣言在成为血肉并像我们一样成为人时，并没有为我们成为大司祭和宗徒，而是另一不同者，即出自妇人之常人；或若有人说祂为自己献祭，而不单为我们献祭——因为那不知罪的祂本不需要献祭——此人当受绝罚。\par

十一、若有人不承认主的肉体是生命之源，并属于天主圣言本身，反而说它属于另一者，虽在尊严上与圣言结合，却只拥有神性的居所；而不如我们所说，承认它是生命之源，因为它属于那赋予万有生命之圣言，此人当受绝罚。\par

十二、若有人不承认天主圣言在肉体中受苦、在肉体中被钉、在肉体中尝了死味，并成为死者中的首生者，因为祂作为天主是生命和生命之源，此人当受绝罚。\par

\SourceRecord{Translated by Blomfield Jackson.；Nicene and Post-Nicene Fathers, Second Series；Vol. 3.；Edited by Philip Schaff and Henry Wace.；Buffalo, NY: Christian Literature Publishing Co.,；1892.；<http://www.newadvent.org/fathers/2701.htm>.}
